\section{Problem Set 2: Determinants}

\begin{enumerate}
\item \textbf{Problem 1}\par
Compute
\[
\det\begin{pmatrix}4&-3\\2&5\end{pmatrix}.
\]

\item \textbf{Problem 2}\par
Find the determinant of
\[
A=\begin{pmatrix}
1&2&0\\
-1&3&4\\
2&0&5
\end{pmatrix}
\]
by expanding along the first row.

\item \textbf{Problem 3}\par
Compute the determinant of
\[
B=\begin{pmatrix}
2&0&0\\
-1&3&0\\
4&5&-2
\end{pmatrix}
\]
without a Laplace expansion. Explain why your method works.

\item \textbf{Problem 4}\par
Expand
\[
\det\begin{pmatrix}
1&0&2\\
3&-1&4\\
0&5&2
\end{pmatrix}
\]
along the row or column that gives the shortest calculation.

\item \textbf{Problem 5}\par
A matrix \(B\) is obtained from a \(3\times3\) matrix \(A\) by interchanging two rows. If \(\det A=-7\), find \(\det B\).

\item \textbf{Problem 6}\par
A matrix \(C\) is obtained from a \(4\times4\) matrix \(A\) by multiplying one row by \(5\). If \(\det A=3\), find \(\det C\).

\item \textbf{Problem 7}\par
A matrix \(D\) is obtained from \(A\) by adding \(4\) times the first row to the third row. If \(\det A=-2\), determine \(\det D\).

\item \textbf{Problem 8}\par
Suppose \(\det A=3\) and \(\det B=-4\) for two \(3\times3\) matrices. Find \(\det(AB)\), \(\det(A^{\mathsf T})\), and \(\det(2A)\).

\item \textbf{Problem 9}\par
Determine whether
\[
A=\begin{pmatrix}
1&2&3\\
2&4&6\\
0&1&1
\end{pmatrix}
\]
is singular without carrying out a full determinant expansion.

\item \textbf{Problem 10}\par
Find all values of \(k\) for which
\[
A(k)=\begin{pmatrix}
k&2\\
3&k-1
\end{pmatrix}
\]
is singular.

\item \textbf{Problem 11}\par
Find \(a\) if
\[
\det\begin{pmatrix}
a&1\\
2&3
\end{pmatrix}=10.
\]

\item \textbf{Problem 12}\par
Compute
\[
\det\begin{pmatrix}
1&2&3\\
0&4&5\\
0&0&6
\end{pmatrix}
\]
and explain the general rule suggested by the result.

\item \textbf{Problem 13}\par
Let
\[
A=\begin{pmatrix}
1&2&0\\
3&1&4\\
2&-1&5
\end{pmatrix}.
\]
Use row operations to transform \(A\) into an upper triangular matrix, keeping track of every change to the determinant. Then compute \(\det A\).

\item \textbf{Problem 14}\par
A \(3\times3\) matrix has determinant \(6\). What happens to its determinant if:
\begin{enumerate}
\item the first and second rows are interchanged;
\item the third row is multiplied by \(-2\);
\item twice the first row is added to the second row?
\end{enumerate}

\item \textbf{Problem 15}\par
The columns of a \(3\times3\) matrix satisfy \(c_3=2c_1-c_2\). What is the determinant? Explain the structural reason.

\item \textbf{Problem 16}\par
The linear transformation
\[
T(x)=
\begin{pmatrix}
2&1\\
0&3
\end{pmatrix}x
\]
acts on the plane. By what factor does \(T\) scale areas? Does it preserve orientation?

\item \textbf{Problem 17}\par
Find the area of the parallelogram spanned by
\[
u=(3,1),
\qquad
v=(-2,4)
\]
using a determinant.

\item \textbf{Problem 18}\par
Compute the Vandermonde determinant
\[
\det\begin{pmatrix}
1&1&1\\
x_1&x_2&x_3\\
x_1^2&x_2^2&x_3^2
\end{pmatrix}
\]
for \(x_1=1\), \(x_2=2\), and \(x_3=4\), first directly and then using the product formula.

\item \textbf{Problem 19}\par
Consider
\[
A(t)=\begin{pmatrix}
1&t&0\\
0&1&t\\
t&0&1
\end{pmatrix}.
\]
Compute \(\det A(t)\) and determine all real values of \(t\) for which the matrix is singular.

\item \textbf{Problem 20}\par
For
\[
A=\begin{pmatrix}
1&2&3\\
2&5&7\\
0&1&4
\end{pmatrix},
\]
compute the determinant in two different ways: by a Laplace expansion and by row reduction. Compare the amount of work and explain which method is more efficient.
\end{enumerate}
